\section{Related Work}
\label{sec:related-work}

\paragraph{Token reduction.}
DynamicViT learns token retention~\citep{dynamicvit2021}, while A-ViT assigns
adaptive depth through halting~\citep{avit2022}. ToMe uses bipartite soft
matching and can compress pretrained or jointly trained transformers
\citep{tome2023}; PiToMe~\citep{pitome2024}, MCTF~\citep{mctf2024}, and
SAD-TM~\citep{sadtm2026} modify its matching signal or merge schedule.
DiffRate treats layerwise budgets as another learned variable
\citep{diffrate2023}. Equal final counts therefore do not imply equal
layerwise computation.

\paragraph{Differentiable merging.}
DTEM is the closest predecessor: it decouples the merge representation from
backbone features and trains it through relaxed grouping~\citep{dtem2024}.
\mn{} builds on its OpenToMe lineage~\citep{opentome}, but places continuous
mass transport and a hard fixed-budget gather in one forward path. The soft
top-$k$ relaxation follows \citet{sujianlin2024topk}; it is not claimed as a
new selection operator. Proportional attention using token size is also
established by ToMe. Our systems addition is the rank-one query/key
reparameterization for retaining that bias in FlashAttention.

\paragraph{Spatial structure and recovery.}
DSM and ALGM combine local and global merging~\citep{dsm2024,algm2024};
GTP-ViT propagates tokens over spatial and semantic edges~\citep{gtpvit2024}.
Recent reducers constrain matching along rows/columns or after space-filling
reordering~\citep{cubistmerge2026,napmat2025}. \mn{} instead constrains
eligible routes by Euclidean radius on the original grid. LookupViT uses
cross-attention between compressed and high-resolution streams
\citep{lookupvit2024}; our recovery read is an architectural component rather
than a claimed general recovery method. Since local merging can reduce
arithmetic without improving throughput~\citep{dsm2024}, we keep analytic
cost and measured runtime distinct.
